\begin{thebibliography}{9}

\bibitem{Ford}
K. Ford, \emph{The distribution of totients},
Ramanujan J. \textbf{2} (1998), 67--151.
The version used here is the revised manuscript
\href{https://arxiv.org/abs/1104.3264v2}{arXiv:1104.3264v2}
(14 July 2013), especially Section~3, Lemma~3.7 and Remark~1.
Author's PDF:
\url{https://www.ford126.web.illinois.edu/wwwpapers/totients.pdf}.

\bibitem{FordLau}
K. Ford and K.-W. Lau, \emph{Problems and Solutions: Solutions:
Asymptotics of a Recurrent Sequence: 10682},
Amer. Math. Monthly \textbf{107}(4) (2000), 374--375.
\href{https://doi.org/10.2307/2589199}{doi:10.2307/2589199}.
Bibliographic identification checked against the University of Illinois
publication record; the full text was not inspected. The year 2000
is used here rather than the year printed in reference [17] of
the revised Ford manuscript.

\bibitem{DLMF}
F. W. J. Olver et al., eds., \emph{NIST Digital Library of Mathematical
Functions}. Section~25.12, especially equation~25.12.12 (convergent
polylogarithm expansion), and Section~5.7 (gamma-function series).
\url{https://dlmf.nist.gov/25.12} and
\url{https://dlmf.nist.gov/5.7}. Accessed 8 October 2026.

\bibitem{Teschl}
G. Teschl, \emph{Topics in Linear and Nonlinear Functional Analysis},
author's textbook page and contents, especially the spectral theorem
and spectral measures in Chapter~5.
\url{https://www.mat.univie.ac.at/~gerald/ftp/book-fa/fa.html}.
The general bounded self-adjoint spectral theorem is the standard
functional-analytic input used in Section~\ref{sec:spectral-notes}.

\bibitem{OpenAIMath}
OpenAI, \emph{An asymptotic formula for the number of totients},
repository manuscript dated 25 September 2026, in \texttt{openai/math}.
Inspected source:
\path{preprints/An-asymptotic-formula-for-the-number-of-totients-September-25-2026/build/source/statements.tex},
commit \path{adc7f1241b42e322a6451854ab7e4b4c146bf78a}.
\url{https://github.com/openai/math}.
Only the explicitly defined recurrence and product scale are used here;
the manuscript's main arithmetic theorem is not an input.

\bibitem{ProveIt}
V. Reshetnikov, \emph{ProveIt}, research and formal-mathematics repository,
snapshot \path{6d78805163d6b5b92d27f8f41555254c49fe0963}.
\url{https://github.com/VladimirReshetnikov/ProveIt}.
The inspected asymptotic and conditioning material was used to select
a nonduplicative direction, not as an unproved lemma in this article.

\end{thebibliography}
